% Generated draft: requires mathematical and source review.
\documentclass[11pt]{article}
\usepackage[margin=25mm]{geometry}
\usepackage{fontspec}
\setmainfont{Latin Modern Roman}
\usepackage{amsmath,amssymb,mathtools,unicode-math}
\setmathfont{Latin Modern Math}
\usepackage{enumitem,needspace,xurl,hyperref}
\setlength{\parindent}{0pt}
\setlength{\parskip}{5pt}
\setlength{\emergencystretch}{3em}
\newcommand{\N}{\mathbb{N}}
\newcommand{\Z}{\mathbb{Z}}
\newcommand{\Q}{\mathbb{Q}}
\newcommand{\R}{\mathbb{R}}
\newcommand{\C}{\mathbb{C}}
\newcommand{\F}{\mathbb{F}}
\newcommand{\A}{\mathbb{A}}
\newcommand{\PP}{\mathbb{P}}
\newcommand{\E}{\mathbb{E}}
\newcommand{\eps}{\varepsilon}
\newcommand{\dd}{\,\mathrm{d}}
\begin{document}
\section*{UnsolvedMath ID: 30001685}
\textbf{Source problem code:} catalogue entry
\subsection{Definitions and mathematical statement}
Work in \(\mathsf{ZFC}\). Put \(\kappa=\aleph_\omega\) and \(\lambda=\kappa^+=\aleph_{\omega+1}\), so \(\operatorname{cf}(\kappa)=\omega\).

A tree is a partially ordered set \((T,<_{T})\) such that the strict predecessors of each node are well ordered by \(<_{T}\). The height of a node is the order type of its predecessors, and \(T_\alpha\) is the set of nodes of height \(\alpha\). A \(\lambda\)-tree has height \(\lambda\) and satisfies \(|T_\alpha|<\lambda\) for every \(\alpha<\lambda\). A cofinal branch is a linearly ordered subset whose node heights are unbounded in \(\lambda\). The tree property at \(\lambda\), denoted here by \(\operatorname{TP}(\lambda)\), means that every \(\lambda\)-tree has a cofinal branch.

A cardinal \(\kappa\) is strong limit if \(2^\mu<\kappa\) for every cardinal \(\mu<\kappa\). The strong-limit formulation of the Singular Cardinal Hypothesis asserts that \(2^\kappa=\kappa^+\) at every singular strong limit cardinal \(\kappa\).

\textbf{Supplied question.} Is it consistent with \(\mathsf{ZFC}\) that the tree property holds at \(\aleph_{\omega+1}\) while the Singular Cardinal Hypothesis fails at \(\aleph_\omega\)?

\textbf{Qualification concerning the intended formulation.} The supplied statement does not explicitly say that \(\aleph_\omega\) is strong limit or define local failure of the Singular Cardinal Hypothesis. Under the strong-limit reading, the precise consistency question is
\[
\operatorname{Con}\!\left(\mathsf{ZFC}+\left[
\operatorname{TP}(\aleph_{\omega+1})\ \land\
(\forall\text{ cardinals }\mu<\aleph_\omega)\;2^\mu<\aleph_\omega\ \land\
2^{\aleph_\omega}>\aleph_{\omega+1}
\right]\right)?
\]
Here \(\operatorname{Con}\) means syntactic consistency, and all cardinal expressions are interpreted within the theory in question. The strong limit condition in this display is an explicit interpretation requiring source confirmation.

Alternatively, a local formulation of SCH at a singular cardinal \(\kappa\) is
\[
2^{\operatorname{cf}(\kappa)}<\kappa
\quad\Longrightarrow\quad
\kappa^{\operatorname{cf}(\kappa)}=\kappa^+.
\]
With this convention, the question instead becomes
\[
\operatorname{Con}\!\left(\mathsf{ZFC}+\left[
\operatorname{TP}(\aleph_{\omega+1})\ \land\
2^{\aleph_0}<\aleph_\omega\ \land\
(\aleph_\omega)^{\aleph_0}>\aleph_{\omega+1}
\right]\right)?
\]
These readings agree when \(\aleph_\omega\) is strong limit; the supplied record does not settle whether that hypothesis is intended. No particular large cardinal assumption or relative consistency strength is specified.
\subsection{Short English statement}
Can a model of set theory have the tree property at \(\aleph_{\omega+1}\) together with failure of the Singular Cardinal Hypothesis at \(\aleph_\omega\)? The source leaves the local meaning of SCH failure implicit, so the draft gives explicit formulations and flags the strong limit hypothesis for review.
\subsection{Sources}
\begin{itemize}
\item Supplied local UnsolvedMath registry record, permanent ID \(30001685\), title: \emph{Tree Property with Failure of Singular Cardinal Hypothesis}; source problem code \texttt{catalogue entry}. The record belongs to the Oberwolfach Reports Open Problems collection and gives proposed year 2011. No individual proposer is supplied.
\item Supplied report citation: \emph{Set Theory} (2011), Oberwolfach Reports, pp. 85--140, DOI 10.4171/owr/2011/02. Source URL and literature-source link labelled ``original OWR report'': \url{https://doi.org/10.4171/owr/2011/02}. The supplied background describes these reports as records of mathematical workshops and their open-problem discussions.
\item Supplied literature-source link labelled ``Sinapova on tree property and SCH'': \url{https://homepages.math.uic.edu/~sinapova/TP%20and%20not%20SCH.pdf}. The supplied assessment attributes to Sinapova an analogous tree-property and SCH-failure model at \(\aleph_{\omega^2}\), and says that this work explicitly leaves the \(\aleph_\omega\) case open. This paraphrases the supplied curation; the result's precise hypotheses and formulation have not been independently verified here.
\item Literature assessment date: \texttt{2026-08-21}. The supplied dated curation labels the problem open and lists the preceding two links as sources checked. Its own qualification requires independent verification of time-sensitive status and cited claims. This draft preserves that dated assessment without independently verifying it or claiming to have visited either source.
\item The supplied registry labels the statement \texttt{corrected\_verified}, with the explanation that TeX notation was enclosed in math delimiters without changing mathematical content. This records a notation correction and does not establish mathematical review of the formulations above.
\end{itemize}
\end{document}
